@@M0@@ is a @@M1@@-ideal.

Now let @@M2@@ be an element of augmentation @@M3@@ of @@M4@@; @@M5@@ is congruent @@M6@@ to @@M7@@, and consequently @@M8@@ is congruent to @@M9@@, whence:
@@M10@@

\sgaStableTarget{sga6.e05.cor.6-6-4}\textbf{Corollary 6.6.4.} \uline{Let} @@M11@@ \uline{be the augmentation homomorphism of} @@M12@@; \uline{then}, \uline{for every} @@M13@@, \uline{we have}:
@@M14@@

By definition, @@M15@@. Since @@M16@@ has augmentation @@M17@@, we can apply Corollary 6.6.2 to it, whence the result.

\sgaStableTarget{sga6.e05.cor.6-6-5}\textbf{Corollary 6.6.5.} \uline{With the notation of} 3.10, @@M18@@.

By the preceding corollary, @@M19@@. The corollary then follows from the fact that the two filtrations are ring filtrations, and that the @@M20@@, for varying @@M21@@ and @@M22@@, generate the @@M23@@-ring filtration.

\sgaStableTarget{sga6.e05.cor.6-6-7}\textbf{Corollary 6.6.7.} \uline{Let} @@M24@@ \uline{be a} \uline{graded} @@M25@@-\uline{algebra of characteristic zero}, \uline{with bounded degrees}; \uline{then the two filtrations defined on} @@M26@@ \uline{are} \uline{identical}.

Suppose that @@M27@@, and let @@M28@@ be an element of @@M29@@. Put @@M30@@; we then have:
@@M31@@, since @@M32@@ has zero augmentation. Consequently,
